\documentclass[12pt, a4paper]{article}
\usepackage[utf8]{inputenc}
\usepackage{amsmath, amssymb, amsthm, amsfonts,amsthm,tikz-cd}
\usepackage{marvosym}
\usepackage{mathrsfs}
\usepackage{CJKutf8}
\usepackage{bm}
\usepackage[margin=1in]{geometry}
\newcommand{\id}{\operatorname{id}}
\newcommand{\qz}{\mathbb{Q}/\mathbb{Z}}
\newcommand{\z}{\mathbb{Z}}
\renewcommand{\hom}{\operatorname{Hom}}
\theoremstyle{remark}
\newtheorem*{remark}{Remark}
\newenvironment{lemma}[1]{
    \noindent \textbf{Lemma #1.} \itshape
}{
    \vspace{1em}
}

% 重新定义 solution 环境，保留标识并在末尾留出 8cm 空白
\newenvironment{solution}
  {\paragraph{\textit{Solution:}}\par\medskip}
  {\hfill$\blacksquare$\vspace{0.1cm}}

\title{Abstract Algebra III: Assignment - Midterm}
\author{
    \begin{CJK*}{UTF8}{gbsn}
    钟皓宇 (12533556)
    \end{CJK*}
}
\date{\today}

\begin{document}

\maketitle


\section*{Lemmas for Problem}
\begin{lemma}{A}
     $\qz$ is an injective $\z$-module.
\end{lemma}
\begin{proof}
    Given an arbitrary $\mathbb{Z}$-module $G$ with its submodule $H$ and a $\mathbb{Z}$-module homomorphism $f:H\to \qz$, let 
    \begin{align*}
        \Omega=\{J\supseteq H\mid \text{ there exists $f_J:J\to \qz$ such that $f_J|_{H}=f$}\}.
    \end{align*}
    According to Zorn's lemma, the partially ordered set $\Omega$ has a maximial element $K$.
    Suppose that $K$ is a proper submodule of $G$.
    Choose an arbitrary element $a\in G\setminus K$, we let $\tilde{K}$ denote the submodule generated by $K$ and $a$.
    
    If the intersection $\langle a \rangle\cap K$ only contains zero, equivalently, $\tilde{K}= K\oplus \mathbb{Z}\cdot a$, then we construct a map $\tilde{f}_K:\tilde{K}\to \qz$  
    by sending $a$ to zero and identified with $f_K$ on $K$. Clearly, it is a $\mathbb{Z}$-module homomorphism since $\tilde{f}_K$ can be identified with the composite of $\tilde{K}\to K \to \qz$ with the first homomorphism is the canonical projection.
    It is impossible because we already set $K$ is the maximial element but we find a larger element $\tilde{K}$ with homomorphism $\tilde{f}_K$.
    
    Next, We let the integer $n$ to be the minimal positive integer that makes $na\in K$ holds. We construct a map $\tilde{f}_K:\tilde{K}\to \qz$ by sending $a$ to $f_K(na)/n$ and $\tilde{f}_K|_K = f_K$. 
    It is a well-defined module homomorphism since if we have $k-ma=0$ in $\tilde{K}$ where $k\in K$ then it implies that $n\mid m $ otherwise there must exists a positive integer $n'<n$ such that $n\equiv n'\bmod m$ and then $a^{n'}\in K$ which is a contradiction.
    Therefore, we have $\tilde{f}_K(k)-\tilde{f}_K(ma)= f_K(k-ma)=0$. 
    
    The process that we extend $f_K$ to $\tilde{f}_K$ implies that $K$ must not be proper, equivalently, $K=G$.
    Then we obtain that $f$ can be extended to $G\to\qz$, it means that $\qz$ is an injective $\mathbb{Z}$-module.
\end{proof}
\begin{lemma}{B}
    Let $A$, $B$, and $C$ be $R$-modules such that $A \subseteq B \subseteq C$.
    If $A$ is an essential submodule of $B$ and $B$ is an essential submodule of $C$, then $A$ is an essential submodule of $C$.
\end{lemma}
\begin{proof}
    We let $\subseteq_e$ denote essential extension.
    Let $W$ be an arbitrary non-zero submodule of $C$. To prove that $A \subseteq_e C$, we need to show that $W \cap A \neq 0$.
    Since $B \subseteq_e C$ and $W$ is a non-zero submodule of $C$, the intersection of $W$ and $B$ must be non-trivial, i.e., $W \cap B \neq 0$. 

    Since $A\subseteq_e B$ and $W \cap B$ is a submodule of $B$, we obtain that 
    \[ (W \cap B) \cap A \neq 0 \]
    which implies that $W\cap A \neq 0$. Therefore, we obtain that $A\subseteq_e C$.
\end{proof}

\begin{lemma}{C}
    Let $A$ and $B$ be $R$-modules such that $A \subseteq_e B$. Let $f:B\to C$ be a $R$-module monomorphism, then $f(A)\subseteq_e f(B)$.
\end{lemma}
\begin{proof}
    Given that an arbitrary submodule $W$ of $f(B)$. Since $A\subseteq_e B$, we have
    \begin{align*}
        f^{-1}(W)\cap A \neq 0.
    \end{align*}
    Since $f$ is monomorphic, we then have 
        \begin{align*}
        W\cap f(A) \neq 0
    \end{align*}
    which implies that $A\subseteq_e B$.
\end{proof}

\section*{Problem 1}

Let $R$ be a ring with $1$.

\begin{itemize}
    \item[(i)] Let $h : U \to V$ be a monomorphism of $R$-modules. Then the following two conditions are equivalent.
    \begin{itemize}
        \item[(a)] A homomorphism $k : V \to Y$ must be a monomorphism whenever the composite map $U \xrightarrow{h} V \xrightarrow{k} Y$ is a monomorphism.
        \item[(b)] If $0 \neq W \subset V$, then $h(U) \cap W \neq 0$.
    \end{itemize}
    The monomorphism $h$ that satisfy the above conditions is said to be \textit{essential}. If there is an injective $R$-module $Q$ with an essential monomorphism $h : V \to Q$, then the diagram $V \xrightarrow{h} Q$ is called an \textit{injective hull} of $V$.

    \item[(ii)] Prove that: any $R$-module $V$ has an injective hull.
\end{itemize}
\begin{solution}
    \textbf{(i)}
    We first show that (a) implies (b). Suppose that $W\neq 0$ and $h(U)\cap W = 0$, then I claim that $U\to V/W$ is injective.
    Since intersection $h(U)\cap W = 0$, we obtain that
    \begin{align*}
        h(U)+W\simeq h(U)\oplus W.
    \end{align*}
    Since $h$ is a monomorphism, there is a canonical isomorphism $h(U)\simeq U$ and a canonical inclusion
    \begin{align*}
        U\hookrightarrow h(U)\oplus W.
    \end{align*}
    In this sense, we find that the homomorphism $U\to V/W$ can be decomposed by 
    \begin{align*}
        U \hookrightarrow h(U)\oplus W \to (h(U)\oplus W)/W.
    \end{align*}
    The last step is exactly the canonical projection, then we obtain that $U\to V/W$ is an injection.
    However, the homomorphism $V\to V/W$ is a monomorphism if and only if $W=(0)$ which contradicts our assumption.

    Conversely, we suppose that $k\circ h$ is a monomorphism but $k$ is not injective. 
    Since $k$ is not injective, it means that $\ker k\neq 0$. Hence, the injectivity of $k\circ h$ implies that $\ker (k\circ h)\supseteq (h(U)\cap \ker k)$ and thus
    $h(U)\cap \ker k=0$. It is a contradiction because we require that for an arbitrary submodule $W$ of $V$ we must have $h(U)\cap W \neq 0$ and $\ker k$ does not satisfy this requirement.

    \textbf{(ii)} For a fixed $R$-module $V$, we construct an injective hull for it as follow.
    According to \textbf{Lemma A.}, we know that $\qz$ is an injective $\z$-module, it equals the functor $\hom_{\z}(-,\qz)$ is exact.
    Let $Q=\hom_\z(R,\qz)$ be the Pontryagin dual module of $R$. Since the tensor-hom duality, we obtain that 
    \begin{align*}
        \hom_R(M,Q)\simeq \hom_\z(M\otimes_R R,\qz)\simeq \hom_\z(M,\qz)
    \end{align*}
    For arbitrary $R$-module $M$.
    It means that the functor $\hom_R(-,Q)$ is also exact, thus $Q$ is an injective $R$-module.
    To construct an injective hull for $R$-module $V$, we first view $V$ as a $\z$-module and consider the $\z$-module homomorphism
    \begin{align*}
        f:V \to \prod_{v\in V}\qz 
    \end{align*} 
    determined bt $\psi_v: V \to \qz$ by extending the $\z$-module homomorphism $\langle v \rangle\to \qz$ by sending $v$ to $1/2+\z$ if $\#\langle v\rangle=\infty$ or to $1/n$ if $\#\langle v\rangle=n$.
    The extension exists because $\qz$ is an injective module. We alse note that $f$ is injective since for an arbitrary element $v\in \ker f$, we have $\psi_v(v) = \z$. It means that the order of $\langle v \rangle$ is one,
    equivalently, $1\cdot v=0$ and thus we have $v=0$. The injectivity of $f$ holds.

    It induces an $R$-module homomorphism 
    \begin{align*}
        \tilde{f}:\hom_\z(R,V) \to \hom_\z(R,\prod_{v\in V}\qz).
    \end{align*}
    Since we have $\hom_\z(R,V)\simeq V$ and $\hom_R(R,\prod_{v\in V}\qz)\simeq \prod_{v\in V} Q$. We can rewrite this homomorphism
    \begin{align*}
        \tilde{f}: V\to \prod_{v\in V} Q.
    \end{align*}

    First of all, the limit $\prod_{v\in V} Q$ is an injective $R$-module since given by an arbitrary $R$-module monomorphism $A \to B$ with homomorphism $A\to \prod_{v\in V} Q$,
    The second homomorphism induces homomorphisms $A\to Q$ for each index $v\in V$ by composite the canonical projection.
    Because $Q$ is an injective $R$-module, there exists lift $B\to Q$ for each $A\to Q$. Hence, by the universal property of production, $B\to Q$ for each index induces a $R$-module homomorphism $B\to \prod_{v\in V} Q$ such that $A\to \prod_{v\in V} Q$
    can be decomposed to $A\to B\to \prod_{v\in V} Q$. It means that $\prod_{v\in V} Q$ is an injective $R$-module.
    Hence, $\tilde{f}$ is injective since $f$ is injective and functor $\hom_\z(R,-)$ is left exact.

    Now, we conisder the partially ordered set
    \begin{align*}
        \Omega = \left\{E\mid E\subseteq \prod_{v\in V} Q \text{ and }F\cap \tilde{f}(V)\neq 0\text{ for any nonzero submodule $F$ of $E$} \right\}.
    \end{align*}
    According to Zorn's lemma, $\Omega$ admits a maximal element we denote it by $E$. Hence, we consider the another partially ordered set 
    \begin{align*}
        \Gamma = \left\{K\mid K\subseteq \prod_{v\in V} Q \text{ and }K\cap E = 0 \right\}.
    \end{align*}
    Similarly, we let $K$ denote a maximal element of $\Gamma$. I claim that $\prod_{v\in V} Q\simeq K\oplus E$.
    To see that we consider a homomorphism $E \to (\prod_{v\in V} Q)/K$. Since $E\cap K =0$, this homomorphism is injective. Therefore, $\prod_{v\in V} Q$ is an injective module implies that 
    there exists a homomorphism 
    \begin{align*}
        \Phi:\left(\prod_{v\in V} Q\right)/K \to \prod_{v\in V} Q
    \end{align*}
    such that the diagram
    \begin{center}
            \begin{tikzcd}[row sep=large, column sep=large]
    % 顶点定义
    E \arrow[r, hook, "\iota"] \arrow[dr,hook, "\psi"'] & 
    \left( \prod_{v \in V} Q \right) \Big/ K \arrow[d, dashed, "\Phi"] \\
    & \prod_{v \in V} Q
    \end{tikzcd}
    \end{center}
    commutes. We notice that $\Phi$ is injective. 
    Because $\psi$ is injective, then $\Phi|_{\iota(E)}$ is injective, therefore we have $\ker \Phi\cap \iota(E) = 0$.
    If $\ker \Phi\neq (0)$,  we consider the canonical quotient 
    \begin{align*}
        \pi:\prod_{v\in V} Q \to \left(\prod_{v\in V} Q\right)/K
    \end{align*}
    and we obtain that $\pi^{-1}(\ker \Phi)\cap E\subseteq K$, since $E\cap K=0$, then we have $\pi^{-1}(\ker \Phi)\cap E=0$.
    Clearly, $\pi^{-1}(\ker \Phi)\not\subseteq K$ then we find that $(K+\pi^{-1}(\ker\Phi))\cap E=0$ and $K\subsetneq(K+\pi^{-1}(\ker\Phi))$.
    It contradicts the maximality of $K$, therefore $\Phi$ must be injective.
    The injectivity of $\Phi$ implies that $\iota(E)$ is an essential submodule of $(\prod_{v\in V} Q)/K$ since if there exists a submodule $M+K$ of $\left(\prod_{v\in V} Q\right)/K$ such that $(M+K)\cap \iota(E)=0$,
    then we have $\Phi(M+K)\cap E = 0$. It implies that $\Phi(M+K)\subseteq K$ equivalently $M+K\in\ker \Phi$.
    Therefore, we obtain that $\iota(E)$ is an essential submodule of $(\prod_{v\in V} Q)/K$.

    According to \textbf{Lemma C.}, the essential extension $\tilde{f}(V)\subseteq_e E$ implies that $\iota\circ\tilde{f}(V)\subseteq_e \iota(E)$ since $\iota$ is injective.
    Hence \textbf{Lemma B.} implies that
    \begin{align*}
        \iota\circ\tilde{f}(V)\subseteq_e (\prod_{v\in V} Q)/K
    \end{align*}
    due to the fact that $\iota\circ\tilde{f}(V)\subseteq_e \iota(E)$ and $\iota(E)\subseteq_e (\prod_{v\in V} Q)/K$.
    The injectivity of $\Phi$ implies that 
    \begin{align*}
        \tilde{f}(V)\subseteq_e \Phi\left((\prod_{v\in V} Q)/K\right).
    \end{align*}
    Since we also have $E\subseteq \Phi\left((\prod_{v\in V} Q)/K\right)$,
    the maximality of $E$ implies that 
    \begin{align*}
        E =  \Phi\left((\prod_{v\in V} Q)/K\right).
    \end{align*}
    Equivalently we say $E \simeq (\prod_{v\in V} Q)/K$ and therefore 
    \begin{align*}
        \prod_{v\in V} Q=K\oplus E.
    \end{align*}
    then $E$ as a summand of injective $R$-module $\prod_{v\in V} Q$ is also a injective $R$-module.
    The definition of $\Omega$ yields $E$ satisfies condition \textbf{(i.b)}, therefore $E$ is an injective hull of $V$.
    Our discussion implies that every $R$-module has an injective hull.
\end{solution} 


\section*{Problem 2}

\begin{itemize}
    \item[(i)] The following holds for an idempotent $e$ of $R$:
    \[ J(eRe) = eJ(R)e = eRe \cap J(R). \]

    \item[(ii)] Prove: $J(M_n(R)) = M_n(J(R))$. Here $M_n(R)$ means the matrix ring.
\end{itemize}
\begin{solution}
    \textbf{(i)} 
    Suppose that $x=er_0e\in J(eRe)$, then for arbitrary $r\in R$, we consider 
    \begin{align*}
        1-rx = (1-e)+(e-rx).
    \end{align*}
    Since $x\in J(eRe)$, then there exixs $t=er_1e\in eRe$ such that 
    \begin{align*}
        t(e-erx) = e.
    \end{align*}
    Hence we have 
    \begin{align*}
        t(e-rx) = er_1e^2(e-rx) = t(e-erx) = e.
    \end{align*}
    Now we consider 
    \begin{align*}
        (1-e+t)(1-rx) = (1-e+t)(1-e+e-rx) = 1-(1-e)rx
    \end{align*}
    Hence we have 
    \begin{align*}
        ((1-e)rx)^2 = (1-e)rx(1-e)rx = 0
    \end{align*}
    since $x\in eRe$ and $x(1-e)=0$.
    Let $\varphi = (1-e+t)$ and $\psi = (1-rx)$ we then have 
    \begin{align*}
        (2\varphi-\varphi\psi\varphi)(1-rx) = 1.
    \end{align*}
    It means that $x\in J(R)$. Equivalently, we prove that $J(eRe)\subseteq J(R)$.

    Since $J(eRe)\subseteq eRe$ we obtain that 
    \begin{align*}
        J(eRe) \subseteq eRe\cap J(R).
    \end{align*}
    Let $x\in eRe\cap J(R)$, then we can suppose that $x=er_0e$ and for any $r\in R$, we have 
    \begin{align*}
        1-rx = 1-r(er_0e)
    \end{align*}
    is invertible. Suppose that $t(1-rx)=1$ with $r=er'e$
    then we consider 
    \begin{align*}
        (ete)(e-rx) = (ete)(e-(er'e)(er_0e)) = et(1-rx)e = e^2 = e.
    \end{align*}
    it implies that 
        \begin{align*}
        J(eRe) \supseteq eRe\cap J(R).
    \end{align*}
    Then we obtain that $J(eRe) = eRe\cap J(R)$. 

    Next, we consider $eJ(R)e$. we consider that $R$ as a left $R$-module $_R R$. Then we have idempotent decomposition 
    \begin{align*}
        _R R = eR\oplus (1-e)R 
    \end{align*}
    Then we have $J(R) = eJ(R)+(1-e)J(R)$ and 
    \begin{align*}
        J(R)\cap eR = eJ(R). \tag{$\star$}
    \end{align*}
    Similarly, if we consider $R$ as a right $R$-module $R_R$, we obtain that 
    \begin{align*}
        J(R)\cap Re = J(R)e. \tag{$\Delta$}
    \end{align*}
    I claim that $eR\cap Re = eRe$. To see that, let $x\in eR\cap Re $ then suppose that $x= er_1=r_2e$. Since $e$ is idempotent, then we have 
    \begin{align*}
        er_1e =xe = r_2e^2=r_2e = x 
    \end{align*}
    it implies that $eR\cap Re\subseteq eRe$. And $eRe \subseteq eR\cap Re$ is easy to see. 
    By the same reason we have $eJ(R)\cap J(R)e = eJ(R)e$. Then we obtain that 
    \begin{align*}
        (J(R)\cap eR)\cap (J(R)\cap Re) = J(R)\cap eRe
    \end{align*}
    and therefore $J(R)\cap eRe = eJ(R)e$.

    \textbf{(ii)} 
    Let $E_{ij}$ denote the matrix of $(i,j)$-entry is $1_R$ and others entires is zero. 
    Let $M$ denote $M_n(R)$.
    Let $A=(a_{ij})$ be an arbitrary element in $J(M)$. Since $J(M)$ is a double sided ideal, then we have $\tilde{A}=E_{1i}AE_{j1}\in J(M)$.
    According to $\textbf{(i)}$, we have 
    \begin{align*}
        E_{11}J(M)E_{11} = J(E_{11}M E_{11}) \simeq J(R).
    \end{align*}
    Then it implies that $E_{11}\tilde{A}E_{11}\in J(E_{11}M E_{11})$. Equivalently, we have $a_{ij}\in J(R)$ and thus $J(M)\subseteq M_{n}(J(R))$.
    Conversely, we first show that $J(R)E_{ii}\subseteq J(M)$ since 
    \begin{align*}
        J(R)E_{ii}= J(E_{ii} M E_{ii})=E_{ii}J(M)E_{ii} \subseteq J(M).
    \end{align*}
    Therefore, we have 
    \begin{align*}
        J(R)E_{ij} = J(R)E_{ii}E_{ij} \subseteq J(M)E_{ij}\subseteq J(M).
    \end{align*}
    It implies that $M_n(J(R))\subseteq J(M)$ and thus 
    \begin{align*}
        J(M_n(R)) = M_n(J(R)).
    \end{align*}
\end{solution}


\section*{Problem 3}
\textbf{Prove:} A submodule of a finitely generated free module over a PID is free.

\begin{solution}
    Let $R$ be an arbitrary PID. We first to show that submodules of $_R R$ is free.
    Let $H$ be an arbitrary submodule of $R$, equivalently, $H$ be an arbitrary ideal of $R$.
    Since $R$ is PID, then there exists an element $x\in R$ such that $H=Rx$. Since $R$ is integeral then there 
    is not nontrivial zero divisor in $R$. Therefore, we find that 
    \begin{align*}
        R \to Rx\quad r\mapsto rx
    \end{align*} 
    is an isomorphism. Equivalently, $H$ is free $R$-module.

    Now, we consider the induction on rank of finitely generated free modules. 
    Suppose that our claim holds for $R^{n-1}$.
    We consider the case that rank of $n$.
    Let $H$ be an arbitrary submodule of $R^n$. Let $p_n$ be the canonical projection on $n^\text{th}$ position of $R^n$.
    Then we obtain a short exact sequence 
    \begin{align*}
        0\to \ker p_n\cap H\to H \to p_n(H) \to 0.
    \end{align*}
    Since $p_n(H)$ is a submodule of $R$, it means that $p_n(H)$ is a free $R$-module. Therefore, the exact sequence split.
    We can write 
    \begin{align*}
        H \simeq (\ker p_n\cap H)\oplus p_n(H).
    \end{align*}
    Since $\ker p_n\cap H$ is the submodule of $R^{n-1}$, then by our assumption, it is also a free $R$-module.
    Therefore, our claim holds.
\end{solution}
\end{document}